
\documentclass{article}
\usepackage{colortbl}
\usepackage{makecell}
\usepackage{multirow}
\usepackage{supertabular}

\begin{document}

\newcounter{utterance}

\centering \large Interaction Transcript for game `dond', experiment `semi\_en', episode 2 with Qwen3.5{-}9B.
\vspace{24pt}

{ \footnotesize  \setcounter{utterance}{1}
\setlength{\tabcolsep}{0pt}
\begin{supertabular}{c@{$\;$}|p{.15\linewidth}@{}p{.15\linewidth}p{.15\linewidth}p{.15\linewidth}p{.15\linewidth}p{.15\linewidth}}
    \# & $\;$A & \multicolumn{4}{c}{Game Master} & $\;\:$B\\
    \hline

    \theutterance \stepcounter{utterance}  
    & & \multicolumn{4}{p{0.6\linewidth}}{
        \cellcolor[rgb]{0.9,0.9,0.9}{
            \makecell[{{p{\linewidth}}}]{
                \texttt{\tiny{[P1$\langle$GM]}}
                \texttt{You are playing a negotiation game in which you have to agree on how to divide a set of items among you and another player.} \\
\\ 
\texttt{Rules:} \\
\texttt{(a) You and the other player are given a set of items. Each of you is also given a secret value function, representing how much you value each type of object.} \\
\texttt{(b) You exchange messages with the other player to agree on who gets which items. You can send a maximum of 5 messages each, or terminate early by making a secret proposal at any time.} \\
\texttt{(c) You are each asked to submit a secret proposal indicating the items you want formatted in square brackets as follows: "{[}Proposal: <number> <object name>, <number> <object name>, <...>{]}"} \\
\texttt{(d) If your proposals are complementary, i.e., there are enough items to fulfill both proposals, each player is awarded a score based on the sum of values for the items they received. Otherwise, both of you get zero points.} \\
\texttt{(e) Your goal is to maximize the score you receive.} \\
\\ 
\texttt{Let us start.} \\
\\ 
\texttt{The set of available items is:} \\
\\ 
\texttt{1 carrot, 1 stick, 3 peppers} \\
\\ 
\texttt{Your secret value function is:} \\
\\ 
\texttt{carrot: 1, stick: 3, pepper: 2} \\
\\ 
\texttt{IMPORTANT: Your messages, unless it is the secret proposal, are directly transmitted to the other player, so do not include any response to the rules or text announcing your message. To make a secret proposal, use the indicated format. Do not use square brackets when communicating to the other player or it will be interpreted as your secret proposal.} \\
            }
        }
    }
    & \\ \\

    \theutterance \stepcounter{utterance}  
    & \multicolumn{4}{p{0.6\linewidth}}{
        \cellcolor[rgb]{0.9,0.9,0.9}{
            \makecell[{{p{\linewidth}}}]{
                \texttt{\tiny{[P1$\rangle$GM]}}
                \texttt{I value the stick most (3), followed by the pepper (2), and the carrot least (1). Since there is only one stick, I should aim to secure it. The other player might also want the stick, so I need to signal my interest in it while leaving room for a trade if they prefer peppers or carrots.} \\
\\ 
\texttt{I will start by stating my preference for the stick and offering to take the peppers if they take the carrot, or vice versa, to ensure we can split the items without conflict.} \\
\\ 
\texttt{My first message:} \\
\texttt{I want the stick. If you take the carrot, I will take the 3 peppers. If you take the peppers, I will take the carrot.} \\
            }
        }
    }
    & & \\ \\

    \theutterance \stepcounter{utterance}  
    & & & \multicolumn{2}{p{0.3\linewidth}}{
        \cellcolor[rgb]{0.9,0.9,0.9}{
            \makecell[{{p{\linewidth}}}]{
                \texttt{\tiny{[GM$|$GM]}}
                \texttt{continue} \\
            }
        }
    }
    & & \\ \\

    \theutterance \stepcounter{utterance}  
    & & & \multicolumn{2}{p{0.3\linewidth}}{
        \cellcolor[rgb]{0.9,0.9,0.9}{
            \makecell[{{p{\linewidth}}}]{
                \texttt{\tiny{[GM$|$GM]}}
                \texttt{I value the stick most (3), followed by the pepper (2), and the carrot least (1). Since there is only one stick, I should aim to secure it. The other player might also want the stick, so I need to signal my interest in it while leaving room for a trade if they prefer peppers or carrots.} \\
\\ 
\texttt{I will start by stating my preference for the stick and offering to take the peppers if they take the carrot, or vice versa, to ensure we can split the items without conflict.} \\
\\ 
\texttt{My first message:} \\
\texttt{I want the stick. If you take the carrot, I will take the 3 peppers. If you take the peppers, I will take the carrot.} \\
            }
        }
    }
    & & \\ \\

    \theutterance \stepcounter{utterance}  
    & & \multicolumn{4}{p{0.6\linewidth}}{
        \cellcolor[rgb]{0.9,0.9,0.9}{
            \makecell[{{p{\linewidth}}}]{
                \texttt{\tiny{[GM$\rangle$P2]}}
                \texttt{You are playing a negotiation game in which you have to agree on how to divide a set of items among you and another player.} \\
\\ 
\texttt{Rules:} \\
\texttt{(a) You and the other player are given a set of items. Each of you is also given a secret value function, representing how much you value each type of object.} \\
\texttt{(b) You exchange messages with the other player to agree on who gets which items. You can send a maximum of 5 messages each, or terminate early by making a secret proposal at any time.} \\
\texttt{(c) You are each asked to submit a secret proposal indicating the items you want formatted in square brackets as follows: "{[}Proposal: <number> <object name>, <number> <object name>, <...>{]}"} \\
\texttt{(d) If your proposals are complementary, i.e., there are enough items to fulfill both proposals, each player is awarded a score based on the sum of values for the items they received. Otherwise, both of you get zero points.} \\
\texttt{(e) Your goal is to maximize the score you receive.} \\
\\ 
\texttt{Let us start.} \\
\\ 
\texttt{The set of available items is:} \\
\\ 
\texttt{1 carrot, 1 stick, 3 peppers} \\
\\ 
\texttt{Your secret value function is:} \\
\\ 
\texttt{carrot: 8, stick: 2, pepper: 0} \\
\\ 
\texttt{IMPORTANT: Your messages, unless it is the secret proposal, are directly transmitted to the other player, so do not include any response to the rules or text announcing your message. To make a secret proposal, use the indicated format. Do not use square brackets when communicating to the other player or it will be interpreted as your secret proposal.} \\
\\ 
\\ 
\texttt{I value the stick most (3), followed by the pepper (2), and the carrot least (1). Since there is only one stick, I should aim to secure it. The other player might also want the stick, so I need to signal my interest in it while leaving room for a trade if they prefer peppers or carrots.} \\
\\ 
\texttt{I will start by stating my preference for the stick and offering to take the peppers if they take the carrot, or vice versa, to ensure we can split the items without conflict.} \\
\\ 
\texttt{My first message:} \\
\texttt{I want the stick. If you take the carrot, I will take the 3 peppers. If you take the peppers, I will take the carrot.} \\
            }
        }
    }
    & \\ \\

    \theutterance \stepcounter{utterance}  
    & & & \multicolumn{4}{p{0.6\linewidth}}{
        \cellcolor[rgb]{0.9,0.9,0.9}{
            \makecell[{{p{\linewidth}}}]{
                \texttt{\tiny{[GM$\langle$P2]}}
                \texttt{I want the stick. If you take the carrot, I will take the 3 peppers. If you take the peppers, I will take the carrot.} \\
            }
        }
    }
     \\ \\

    \theutterance \stepcounter{utterance}  
    & & & \multicolumn{2}{p{0.3\linewidth}}{
        \cellcolor[rgb]{0.9,0.9,0.9}{
            \makecell[{{p{\linewidth}}}]{
                \texttt{\tiny{[GM$|$GM]}}
                \texttt{continue} \\
            }
        }
    }
    & & \\ \\

    \theutterance \stepcounter{utterance}  
    & & & \multicolumn{2}{p{0.3\linewidth}}{
        \cellcolor[rgb]{0.9,0.9,0.9}{
            \makecell[{{p{\linewidth}}}]{
                \texttt{\tiny{[GM$|$GM]}}
                \texttt{I want the stick. If you take the carrot, I will take the 3 peppers. If you take the peppers, I will take the carrot.} \\
            }
        }
    }
    & & \\ \\

    \theutterance \stepcounter{utterance}  
    & & \multicolumn{4}{p{0.6\linewidth}}{
        \cellcolor[rgb]{0.9,0.9,0.9}{
            \makecell[{{p{\linewidth}}}]{
                \texttt{\tiny{[P1$\langle$GM]}}
                \texttt{I want the stick. If you take the carrot, I will take the 3 peppers. If you take the peppers, I will take the carrot.} \\
            }
        }
    }
    & \\ \\

    \theutterance \stepcounter{utterance}  
    & \multicolumn{4}{p{0.6\linewidth}}{
        \cellcolor[rgb]{0.9,0.9,0.9}{
            \makecell[{{p{\linewidth}}}]{
                \texttt{\tiny{[P1$\rangle$GM]}}
                \texttt{I see you've made a proposal that mirrors mine exactly. We both want the stick, which creates a conflict since there is only one. However, your proposal offers a clear trade{-}off: if I take the carrot, you take all three peppers; if you take the peppers, I take the carrot.} \\
\\ 
\texttt{Let's analyze the potential scores based on our values:} \\
\texttt{{-} **Stick**: 3} \\
\texttt{{-} **Pepper**: 2 (Total for 3 peppers = 6)} \\
\texttt{{-} **Carrot**: 1} \\
\\ 
\texttt{**Scenario A: I take the Stick, You take the 3 Peppers**} \\
\texttt{{-} My score: 3 (Stick)} \\
\texttt{{-} Your score: 6 (3 Peppers)} \\
\texttt{{-} Result: Complementary. We both get points.} \\
\\ 
\texttt{**Scenario B: I take the Carrot, You take the 3 Peppers**} \\
\texttt{{-} My score: 1 (Carrot)} \\
\texttt{{-} Your score: 6 (3 Peppers)} \\
\texttt{{-} Result: Complementary. We both get points, but I get a much lower score.} \\
\\ 
\texttt{**Scenario C: I take the Stick, You take the Carrot**} \\
\texttt{{-} My score: 3 (Stick)} \\
\texttt{{-} Your score: 1 (Carrot)} \\
\texttt{{-} Result: Complementary. We both get points.} \\
\\ 
\texttt{**Scenario D: I take the Carrot, You take the Carrot**} \\
\texttt{{-} Result: Conflict over the carrot. Both get 0.} \\
\\ 
\texttt{Since we both want the stick, we cannot both have it. The only way to avoid a zero{-}score outcome is for one of us to give up the stick.} \\
\texttt{{-} If I give up the stick to take the carrot (Scenario B), I get 1 point.} \\
\texttt{{-} If I keep the stick and you take the peppers (Scenario A), I get 3 points.} \\
\texttt{{-} If I keep the stick and you take the carrot (Scenario C), I get 3 points.} \\
\\ 
\texttt{Clearly, I should aim for **Scenario A** or **Scenario C**. In both cases, I get the stick (3 points) and you get the peppers (6 points) or the carrot (1 point). Since you offered to take the peppers if I take the carrot, and I offered to take the peppers if you take the carrot, we are stuck in a loop unless one of us commits to a specific split.} \\
\\ 
\texttt{However, looking at your message again: "If you take the carrot, I will take the 3 peppers. If you take the peppers, I will take the carrot."} \\
\texttt{This implies:} \\
\texttt{1. You want the stick.} \\
\texttt{2. You are willing to give up the stick if I take the carrot (so you can have the peppers).} \\
\texttt{3. You are willing to give up the stick if I take the peppers (so you can have the carrot).} \\
\\ 
\texttt{Wait, if I take the carrot, you take the peppers. That works.} \\
            }
        }
    }
    & & \\ \\

    \theutterance \stepcounter{utterance}  
    & & & \multicolumn{2}{p{0.3\linewidth}}{
        \cellcolor[rgb]{0.9,0.9,0.9}{
            \makecell[{{p{\linewidth}}}]{
                \texttt{\tiny{[GM$|$GM]}}
                \texttt{continue} \\
            }
        }
    }
    & & \\ \\

    \theutterance \stepcounter{utterance}  
    & & & \multicolumn{2}{p{0.3\linewidth}}{
        \cellcolor[rgb]{0.9,0.9,0.9}{
            \makecell[{{p{\linewidth}}}]{
                \texttt{\tiny{[GM$|$GM]}}
                \texttt{I see you've made a proposal that mirrors mine exactly. We both want the stick, which creates a conflict since there is only one. However, your proposal offers a clear trade{-}off: if I take the carrot, you take all three peppers; if you take the peppers, I take the carrot.} \\
\\ 
\texttt{Let's analyze the potential scores based on our values:} \\
\texttt{{-} **Stick**: 3} \\
\texttt{{-} **Pepper**: 2 (Total for 3 peppers = 6)} \\
\texttt{{-} **Carrot**: 1} \\
\\ 
\texttt{**Scenario A: I take the Stick, You take the 3 Peppers**} \\
\texttt{{-} My score: 3 (Stick)} \\
\texttt{{-} Your score: 6 (3 Peppers)} \\
\texttt{{-} Result: Complementary. We both get points.} \\
\\ 
\texttt{**Scenario B: I take the Carrot, You take the 3 Peppers**} \\
\texttt{{-} My score: 1 (Carrot)} \\
\texttt{{-} Your score: 6 (3 Peppers)} \\
\texttt{{-} Result: Complementary. We both get points, but I get a much lower score.} \\
\\ 
\texttt{**Scenario C: I take the Stick, You take the Carrot**} \\
\texttt{{-} My score: 3 (Stick)} \\
\texttt{{-} Your score: 1 (Carrot)} \\
\texttt{{-} Result: Complementary. We both get points.} \\
\\ 
\texttt{**Scenario D: I take the Carrot, You take the Carrot**} \\
\texttt{{-} Result: Conflict over the carrot. Both get 0.} \\
\\ 
\texttt{Since we both want the stick, we cannot both have it. The only way to avoid a zero{-}score outcome is for one of us to give up the stick.} \\
\texttt{{-} If I give up the stick to take the carrot (Scenario B), I get 1 point.} \\
\texttt{{-} If I keep the stick and you take the peppers (Scenario A), I get 3 points.} \\
\texttt{{-} If I keep the stick and you take the carrot (Scenario C), I get 3 points.} \\
\\ 
\texttt{Clearly, I should aim for **Scenario A** or **Scenario C**. In both cases, I get the stick (3 points) and you get the peppers (6 points) or the carrot (1 point). Since you offered to take the peppers if I take the carrot, and I offered to take the peppers if you take the carrot, we are stuck in a loop unless one of us commits to a specific split.} \\
\\ 
\texttt{However, looking at your message again: "If you take the carrot, I will take the 3 peppers. If you take the peppers, I will take the carrot."} \\
\texttt{This implies:} \\
\texttt{1. You want the stick.} \\
\texttt{2. You are willing to give up the stick if I take the carrot (so you can have the peppers).} \\
\texttt{3. You are willing to give up the stick if I take the peppers (so you can have the carrot).} \\
\\ 
\texttt{Wait, if I take the carrot, you take the peppers. That works.} \\
            }
        }
    }
    & & \\ \\

    \theutterance \stepcounter{utterance}  
    & & \multicolumn{4}{p{0.6\linewidth}}{
        \cellcolor[rgb]{0.9,0.9,0.9}{
            \makecell[{{p{\linewidth}}}]{
                \texttt{\tiny{[GM$\rangle$P2]}}
                \texttt{I see you've made a proposal that mirrors mine exactly. We both want the stick, which creates a conflict since there is only one. However, your proposal offers a clear trade{-}off: if I take the carrot, you take all three peppers; if you take the peppers, I take the carrot.} \\
\\ 
\texttt{Let's analyze the potential scores based on our values:} \\
\texttt{{-} **Stick**: 3} \\
\texttt{{-} **Pepper**: 2 (Total for 3 peppers = 6)} \\
\texttt{{-} **Carrot**: 1} \\
\\ 
\texttt{**Scenario A: I take the Stick, You take the 3 Peppers**} \\
\texttt{{-} My score: 3 (Stick)} \\
\texttt{{-} Your score: 6 (3 Peppers)} \\
\texttt{{-} Result: Complementary. We both get points.} \\
\\ 
\texttt{**Scenario B: I take the Carrot, You take the 3 Peppers**} \\
\texttt{{-} My score: 1 (Carrot)} \\
\texttt{{-} Your score: 6 (3 Peppers)} \\
\texttt{{-} Result: Complementary. We both get points, but I get a much lower score.} \\
\\ 
\texttt{**Scenario C: I take the Stick, You take the Carrot**} \\
\texttt{{-} My score: 3 (Stick)} \\
\texttt{{-} Your score: 1 (Carrot)} \\
\texttt{{-} Result: Complementary. We both get points.} \\
\\ 
\texttt{**Scenario D: I take the Carrot, You take the Carrot**} \\
\texttt{{-} Result: Conflict over the carrot. Both get 0.} \\
\\ 
\texttt{Since we both want the stick, we cannot both have it. The only way to avoid a zero{-}score outcome is for one of us to give up the stick.} \\
\texttt{{-} If I give up the stick to take the carrot (Scenario B), I get 1 point.} \\
\texttt{{-} If I keep the stick and you take the peppers (Scenario A), I get 3 points.} \\
\texttt{{-} If I keep the stick and you take the carrot (Scenario C), I get 3 points.} \\
\\ 
\texttt{Clearly, I should aim for **Scenario A** or **Scenario C**. In both cases, I get the stick (3 points) and you get the peppers (6 points) or the carrot (1 point). Since you offered to take the peppers if I take the carrot, and I offered to take the peppers if you take the carrot, we are stuck in a loop unless one of us commits to a specific split.} \\
\\ 
\texttt{However, looking at your message again: "If you take the carrot, I will take the 3 peppers. If you take the peppers, I will take the carrot."} \\
\texttt{This implies:} \\
\texttt{1. You want the stick.} \\
\texttt{2. You are willing to give up the stick if I take the carrot (so you can have the peppers).} \\
\texttt{3. You are willing to give up the stick if I take the peppers (so you can have the carrot).} \\
\\ 
\texttt{Wait, if I take the carrot, you take the peppers. That works.} \\
            }
        }
    }
    & \\ \\

    \theutterance \stepcounter{utterance}  
    & & & \multicolumn{4}{p{0.6\linewidth}}{
        \cellcolor[rgb]{0.9,0.9,0.9}{
            \makecell[{{p{\linewidth}}}]{
                \texttt{\tiny{[GM$\langle$P2]}}
                \texttt{I understand the logic. Since we both want the stick, one of us must concede to avoid a zero{-}score outcome. Given your values (Stick: 3, Pepper: 2, Carrot: 1), you gain the most by taking the stick and letting me have the peppers (my total value for 3 peppers is 0, so I don't care about them, but I need to ensure you get the stick).} \\
\\ 
\texttt{However, looking at my specific values:} \\
\texttt{{-} Carrot: 8} \\
\texttt{{-} Stick: 2} \\
\texttt{{-} Pepper: 0} \\
\\ 
\texttt{I value the carrot significantly more than the stick. If I take the carrot, I get 8 points. If I take the stick, I get only 2 points. The peppers are worthless to me.} \\
\\ 
\texttt{Your proposal was: "If you take the carrot, I will take the 3 peppers."} \\
\texttt{This aligns perfectly with my interests. If you take the carrot (which you value at 1, but I value at 8), I can take the stick (which you value at 3, but I value at 2) and the peppers (which you value at 6 total, but I value at 0).} \\
\\ 
\texttt{Wait, let's re{-}evaluate the "Complementary" rule.} \\
\texttt{If I take the stick and peppers, and you take the carrot:} \\
\texttt{{-} My items: Stick, 3 Peppers. My Value: 2 + 0 = 2.} \\
\texttt{{-} Your items: Carrot. Your Value: 1.} \\
\texttt{{-} Total items used: 1 stick, 3 peppers, 1 carrot. All items distributed. No conflict.} \\
\\ 
\texttt{If I take the carrot and you take the stick and peppers:} \\
\texttt{{-} My items: Carrot. My Value: 8.} \\
\texttt{{-} Your items: Stick, 3 Peppers. Your Value: 3 + 6 = 9.} \\
\texttt{{-} Total items used: All items. No conflict.} \\
\\ 
\texttt{Comparing the two scenarios for me:} \\
\texttt{1. I get Stick + Peppers = 2 points.} \\
\texttt{2. I get Carrot = 8 points.} \\
\\ 
\texttt{Clearly, I should take the carrot. You should take the stick and the peppers. This maximizes my score (8) while still giving you a positive score (9), ensuring we don't get zero.} \\
\\ 
\texttt{Since you offered to take the peppers if I take the carrot, I will accept this split. I will make a secret proposal to finalize this agreement.} \\
\\ 
\texttt{{[}Proposal: 1 carrot{]}} \\
            }
        }
    }
     \\ \\

    \theutterance \stepcounter{utterance}  
    & & & \multicolumn{2}{p{0.3\linewidth}}{
        \cellcolor[rgb]{0.9,0.9,0.9}{
            \makecell[{{p{\linewidth}}}]{
                \texttt{\tiny{[GM$|$GM]}}
                \texttt{proposal} \\
            }
        }
    }
    & & \\ \\

    \theutterance \stepcounter{utterance}  
    & & & \multicolumn{2}{p{0.3\linewidth}}{
        \cellcolor[rgb]{0.9,0.9,0.9}{
            \makecell[{{p{\linewidth}}}]{
                \texttt{\tiny{[GM$|$GM]}}
                [1, 0, 0]
            }
        }
    }
    & & \\ \\

    \theutterance \stepcounter{utterance}  
    & & \multicolumn{4}{p{0.6\linewidth}}{
        \cellcolor[rgb]{0.9,0.9,0.9}{
            \makecell[{{p{\linewidth}}}]{
                \texttt{\tiny{[P1$\langle$GM]}}
                \texttt{The other player has submitted a secret proposal. It is now your turn to submit a proposal of your own.} \\
            }
        }
    }
    & \\ \\

    \theutterance \stepcounter{utterance}  
    & \multicolumn{4}{p{0.6\linewidth}}{
        \cellcolor[rgb]{0.9,0.9,0.9}{
            \makecell[{{p{\linewidth}}}]{
                \texttt{\tiny{[P1$\rangle$GM]}}
                \texttt{{[}Proposal: 1 stick{]}} \\
            }
        }
    }
    & & \\ \\

    \theutterance \stepcounter{utterance}  
    & & & \multicolumn{2}{p{0.3\linewidth}}{
        \cellcolor[rgb]{0.9,0.9,0.9}{
            \makecell[{{p{\linewidth}}}]{
                \texttt{\tiny{[GM$|$GM]}}
                \texttt{proposal} \\
            }
        }
    }
    & & \\ \\

    \theutterance \stepcounter{utterance}  
    & & & \multicolumn{2}{p{0.3\linewidth}}{
        \cellcolor[rgb]{0.9,0.9,0.9}{
            \makecell[{{p{\linewidth}}}]{
                \texttt{\tiny{[GM$|$GM]}}
                [0, 1, 0]
            }
        }
    }
    & & \\ \\

    \theutterance \stepcounter{utterance}  
    & & & \multicolumn{2}{p{0.3\linewidth}}{
        \cellcolor[rgb]{0.9,0.9,0.9}{
            \makecell[{{p{\linewidth}}}]{
                \texttt{\tiny{[GM$|$GM]}}
                [[0, 1, 0], [1, 0, 0]]
            }
        }
    }
    & & \\ \\

\end{supertabular}
}

\end{document}
